\section*{Discussion}

The measurements confirm a single mechanism, and the theory names it. Parity labels make a symmetry-mandated zero structural: Theorem~\ref{thm:zero} holds at any parameter values, so the zero cannot be trained away because it was never trained in, and Corollary~\ref{cor:jac} makes it differentiate correctly. Removing the labels converts a certainty into an ordinary regression output, which on a centrosymmetric crystal is an ordinary regression error the size of the property itself. The dipole control shows the failure is not a property of odd-parity targets as such but of evaluating them where symmetry makes the answer exact; the point-group split shows SO(3) models retain precisely the rotational part of Neumann's principle and nothing more, up to the ceiling Corollary~\ref{cor:rot} fixes in advance; and the agreement between architecturally unrelated models places the effect in the symmetry class rather than in any network.

This bears directly on whether symmetry can be learned from data rather than built in. On rotations the evidence for that proposition is real: augmentation closes the data-efficiency gap given enough epochs, and unconstrained potentials trained on large corpora are competitive on accuracy and speed. Our augmentation experiment poses the same proposition where the symmetry is exact, and it fails. A thousand centrosymmetric crystals carrying the exact answer, with their loss weight raised a hundredfold, move the false-flag fraction only from 0.90 to 0.68 on the very crystals the model was shown, and leave the violation median more than four orders of magnitude above the O(3) floor. Gradient descent moves the predictions towards zero and does not arrive, because the target is not a value to be fitted but an identity to be satisfied. The distinction that matters is therefore not equivariant versus unconstrained, nor how much data is available, but whether the quantity predicted is an approximation or an exactness. Where the physics permits any value, learned symmetry is a reasonable engineering trade; where it permits exactly one, it is no trade at all, because an approximate constraint delivers an error the size of the property. Not all symmetry, then, can be learned from data. \MISSING{the title now makes a learnability claim and the evidence for it is one augmentation size (1,000 zero-labelled crystals) plus a loss-weight sweep on one architecture. Sweep the augmentation size over N = 0, 250, 1,000, 4,000 and 16,000 zero-labelled crystals, three seeds each, and plot the false-flag fraction and violation median against N. If the curve plateaus well above zero, as the theory predicts, this becomes the most-cited panel in the paper. If it falls towards zero, change the title. See EXPERIMENTS\_TO\_RUN.md, B-3} On the evidence here the dividing line is exactness: a constraint that permits a range of values can be approached by fitting, while one that permits a single value is not attained by it, even when that value is placed in the training set and weighted a hundredfold. Non-conservative force models made the same bargain against energy conservation, and when they were assessed against the physics rather than against the test error, they were found wanting\cite{bigi2025forces}.

The result also gives quantitative content to a distinction the field treats as terminology. Curie's principle, that an equivariant map cannot produce outputs less symmetric than its inputs, is the foundation on which equivariant networks are advertised. For improper operations it is silently violated by a widely deployed subclass of the very models it is invoked to describe: upholding it for tensor properties requires the full orthogonal group, and rotational equivariance upholds only its proper half.

The distinction matters for how symmetry-correct predictions should be engineered. Space-group-informed modules, canonicalization and symmetry-adapted output decompositions repair the output of a model whose features remain parity-blind, and Corollary~\ref{cor:sym} makes the cost concrete: such a repair works only for the quantity someone remembered to repair, requires the target's parity to be known in advance, and leaves everything else derived from the same features unprotected. Feature-level parity makes the guarantee available to every quantity the features feed, derivatives included, provided the readout declares the target's parity correctly. This is not an argument against readout-level methods, which have accuracy benefits of their own, but an argument that the feature symmetry group is a prior design decision, made once for the whole model and currently made invisibly. As property heads are attached to the embeddings of pretrained universal potentials, that decision is inherited by every downstream task: a parity-even rank-2 target such as the dielectric tensor never exposes it, and the first parity-odd head attached to a rotation-only backbone will fail in the way measured here. \MISSING{this is the paper's most consequential practical claim and it is not tested: every model here is trained from scratch. Fine-tune a piezoelectric head on two pretrained universal potentials, one parity-aware (MACE-MP-0) and one parity-blind (eSEN or UMA), with the backbone frozen, and report the false-flag fractions. That is the scenario practitioners are actually in. See EXPERIMENTS\_TO\_RUN.md, H-3}

The practical rule is narrow enough to check before a model is built. Parity labels are required whenever the target is parity-odd and the evaluation population may contain structures whose point group forbids it. Corollary~\ref{cor:rot} makes the second condition precise: rotation alone suffices only where the proper-rotation subgroup already annihilates the target, which for a rank-3 tensor means class m$\bar 3$m and nothing else, 8.3\% of the crystals evaluated here. The rule is not confined to piezoelectricity: in the electric-dipole approximation the second-harmonic-generation tensor is also rank-3 and parity-odd and vanishes identically in every centrosymmetric crystal, and it is an active target for machine-learned discovery of nonlinear-optical materials\cite{optixnet2026}, where a rotation-only backbone is unconstrained for the same reason.

One objection deserves a direct answer: a screening pipeline can filter centrosymmetric candidates out before the model sees them, as high-throughput searches for nonlinear-optical crystals routinely do. Pre-filtering is the same repair with the same costs. It protects only the quantity someone remembered to filter for, it requires the crystal's symmetry to be known before the prediction rather than after, and it is only as reliable as the symmetry determination feeding it. That last cost is measurable in our own data: all 2,000 evaluation crystals are centrosymmetric at the selection tolerance, but 2.2\% are not at $10^{-4}$ and 6.5\% are not at $10^{-5}$, and structures from relaxation or from generative models arrive with no reliable space-group label at all. A parity-labelled model needs no filter and behaves correctly at the boundary: the single raw-coordinate structure an O(3) arm flagged proved genuinely non-centrosymmetric below the selection tolerance, so the model had responded to real symmetry breaking in its input.

Sensitive resonance techniques do detect piezoelectric responses in nominally centrosymmetric materials\cite{aktas2021nominally}, including NaCl and CaF$_2$, which appear in our named-materials table. Those measurements report a real sample whose actual structure breaks inversion, whether or not an extrinsic cause is apparent, and their calibrated effective coefficients are far below conventional piezoelectrics: Aktas et al.\ report the symmetry-forbidden effect spanning roughly $1$\,pm\,V$^{-1}$ down to $10^{-5}$\,pm\,V$^{-1}$, about $0.5\%$ to $2\times10^{-7}\%$ of poled lead zirconate titanate, with the low-defect ionic crystals NaCl and CaF$_2$ at the bottom of that range, below the detection limit of conventional piezoelectric measurement; the SO(3) violations here occur on exactly symmetric digital structures and reach the full magnitude of genuine piezoelectric tensors. The parity-labelled models capture the legitimate version of the phenomenon, as the raw-coordinate case above shows.

EquiformerV2's behaviour is reported strictly as measured properties of its released code, reproduced against the pinned upstream source: it is nondeterministic in evaluation because it draws a random per-edge reference frame on every pass, its rotational equivariance is approximate at the ten-percent level, and its autograd Jacobian is incomplete (Supplementary Note~11 and Supplementary Tables~16--17). No conclusion depends on this, and one part of it is informative: the stability clause of Corollary~\ref{cor:rot} predicts that a model with defects of that size need not respect the rotation ceiling, and EquiformerV2 does not. Its inclusion serves one purpose, to show the headline behaviour in a deployed model our construction never touched.

The study has limits. All three matched O(3) arms realize their parity labels through the same library, e3nn\cite{geiger2022e3nn}; an intended non-e3nn control was withdrawn for numerical conditioning reasons unrelated to parity, documented in Supplementary Note~12, so implementation-specific contributions to the O(3) side cannot be fully excluded. Three seeds limit seed-level statistics, and all statistical weight comes from the structure-level paired tests. The violation magnitude is extensive, though a size-normalized variant leaves every conclusion unchanged (Supplementary Note~4). The scalar-energy control is deliberately unconverged. The augmentation experiment used one budget and the loss-weight sweep one architecture; neither excludes that more aggressive regimes push violations lower, but the in-training control bounds what any regime can achieve. Hyperparameters were shared across cores rather than tuned. Finally, the prevalence audit is pinned to the versions listed in Supplementary Table~1 and does not by itself cover models released since \MISSING{either add eSEN, UMA, and EquiformerV3 to the audit under the same protocol, or state here that they are untested; see the red flag in Supplementary Note~2}.

Whether a machine-learned model can make physically impossible tensor predictions is therefore decided before training begins, by whether its features carry parity labels, a property nine of the fifteen architectures we audited lack and none advertises. Because it is not advertised it cannot be checked: a model card that reports a maximum angular degree but not a feature symmetry group has withheld the fact that decides whether the model can respect the improper half of Neumann's principle. For parity-odd targets, and for any pipeline that may one day attach one, the label is not bookkeeping; it is the physics.
